\documentclass[11pt,letterpaper]{article}
\usepackage[top=0.7in,bottom=0.7in,left=0.75in,right=0.75in]{geometry}
\usepackage[T1]{fontenc}
\usepackage[utf8]{inputenc}
\usepackage{lmodern}
\usepackage{microtype}
\usepackage{enumitem}
\usepackage[hidelinks]{hyperref}
\usepackage{titlesec}
\setlength{\parindent}{0pt}
\setlength{\parskip}{0pt}
\pagestyle{empty}
\titleformat{\section}{\normalfont\bfseries}{}{0em}{\uppercase}[\vspace{-4pt}\hrule\vspace{4pt}]
\titlespacing{\section}{0pt}{8pt}{5pt}
\setlist[itemize]{leftmargin=1.2em,label=--,itemsep=0pt,topsep=2pt,parsep=0pt}
\begin{document}
\begin{center}
{\fontsize{16}{19}\selectfont\textbf{\MakeUppercase{Diego Castillo Pineda}}}\\[4pt]
{\small La Florida, Santiago, Chile \quad|\quad +56 9 94162547 \quad|\quad \href{mailto:diego.castillop11@gmail.com}{diego.castillop11@gmail.com}\\
\href{https://https://www.linkedin.com/in/diego-castillo-pineda/}{https://www.linkedin.com/in/diego-castillo-pineda/} \quad|\quad \href{https://https://github.com/diegocastillop11-cloud}{https://github.com/diegocastillop11-cloud}}
\end{center}

\section{Resumen Profesional}
Analista de Datos con experiencia sólida en SQL Server, Python y automatización de procesos en entornos de alta concurrencia para industrias financieras y de entretenimiento masivo. Especializado en optimización de bases de datos, reportería avanzada con Power BI y desarrollo de soluciones analíticas en Microsoft Azure. Reduje errores operativos en un 40\% mediante automatización T-SQL en Punto Ticket, gestionando infraestructura crítica para eventos con decenas de miles de transacciones simultáneas. Busco aportar mi experiencia técnica en análisis de datos y diseño de soluciones escalables como Ssr. Data Analyst en Coderio para el sector financiero.

\section{Experiencia Profesional}

\textbf{Punto Ticket} \hfill Santiago, Chile\\
\textit{Analista de Base de Datos Senior} \hfill Nov. 2019 – Feb. 2026
\begin{itemize}
  \item Diseñé y optimicé stored procedures en SQL Server para eventos masivos, procesando más de 50,000 transacciones por hora con tiempos de respuesta bajo 200ms, garantizando estabilidad operativa en períodos críticos de venta.
  \item Automaticé procesos de carga de datos mediante scripts T-SQL y Python, reduciendo en 40\% los errores manuales y disminuyendo el tiempo de ejecución de procesos recurrentes de 4 horas a 45 minutos.
  \item Desarrollé consultas dinámicas con PIVOT y análisis multidimensional para generar reportería ejecutiva en Power BI, reemplazando procesos manuales de Excel y entregando insights operativos en tiempo real a stakeholders clave.
  \item Integré sistemas heterogéneos mediante APIs REST utilizando JSON y XML, facilitando la interoperabilidad entre plataformas de ticketing, CRM y sistemas financieros.
  \item Participé en la migración de bases de datos on-premise a Microsoft Azure SQL Database, colaborando en la estrategia de cloud adoption y optimización de costos operativos.
  \item Administré bases de datos SQL Server en entornos de alta disponibilidad, implementando respaldos automatizados, monitoreo de performance y planes de recuperación ante desastres.
\end{itemize}
\vspace{2pt}
\textbf{Imperial S.A.} \hfill Santiago, Chile\\
\textit{Analista Funcional} \hfill Jun. 2017 – Nov. 2019
\begin{itemize}
  \item Analicé requerimientos funcionales y traduje necesidades de negocio en especificaciones técnicas para equipos de desarrollo, trabajando con SQL Server, Oracle y sistemas ERP en ambientes productivos.
  \item Desarrollé consultas SQL avanzadas y reportes personalizados para áreas financieras y operativas, optimizando la extracción de datos críticos para la toma de decisiones ejecutivas.
  \item Realicé parametrización y configuración de sistemas ERP (Oracle, SAP), soportando procesos de contabilidad, logística e inventario mediante levantamiento de requerimientos con usuarios clave.
  \item Documenté flujos de proceso, instructivos técnicos y manuales de usuario, facilitando la capacitación de equipos y la estandarización de procedimientos operativos.
  \item Brindé soporte funcional y resolución de incidencias en ambientes productivos, coordinando con equipos técnicos la corrección de errores y mejoras en aplicaciones internas.
\end{itemize}
\vspace{2pt}
\textbf{OB GROUP PARK} \hfill Santiago, Chile\\
\textit{Analista de Sistemas y TI} \hfill Mayo 2014 – Abril 2017
\begin{itemize}
  \item Lideré proyectos de integración tecnológica para procesos contables y de gestión en retail, coordinando con equipos multidisciplinarios y asegurando la continuidad operativa durante implementaciones críticas.
  \item Implementé soluciones de bases de datos SQL Server para mejorar la trazabilidad de transacciones financieras, reduciendo el tiempo de cierre contable mensual en 25\%.
  \item Colaboré con áreas de negocio para la adopción de nuevas tecnologías, capacitando usuarios finales y documentando procedimientos operativos estándar.
\end{itemize}
\vspace{2pt}
\textbf{Hewlett-Packard Company} \hfill Santiago, Chile\\
\textit{Agente de Mesa de Ayuda Nivel 1} \hfill Ago. 2013 – Mayo 2014
\begin{itemize}
  \item Brindé soporte técnico remoto a usuarios finales, diagnosticando y resolviendo incidencias de hardware, software y conectividad con un SLA de resolución del 95\% en primer contacto.
\end{itemize}

\section{Proyectos Destacados}

\textbf{Sistema de Análisis Financiero con IA} \hfill 2024
\begin{itemize}
  \item Desarrollé una aplicación full-stack con React y Node.js que utiliza prompt engineering con Claude AI para generar insights financieros automáticos desde datasets SQL Server.
  \item Implementé base de datos PostgreSQL segura con autenticación JWT y encriptación de datos sensibles, aplicando principios de UI/UX para dashboards interactivos con visualizaciones en tiempo real.
  \item Automaticé procesos de ETL con Python y Docker, reduciendo el tiempo de generación de reportes financieros de 2 horas a 10 minutos mediante pipelines de datos escalables.
\end{itemize}

\section{Habilidades T\'{e}cnicas}
\textbf{Análisis de Datos \& Bases de Datos:} SQL Server (Experto), T-SQL, Azure SQL Database, Oracle, PostgreSQL, MongoDB, Optimización de Queries \\
\textbf{Programación \& IA:} Python (Pandas, NumPy), C\#, JavaScript (ES6+), Node.js, Prompt Engineering (Claude, Gemini), RESTful APIs \\
\textbf{Visualización \& Reporting:} Power BI (DAX, Power Query), Excel Avanzado (VBA, Tablas Dinámicas), Business Analytics, Dashboards Interactivos \\
\textbf{Cloud \& DevOps:} Microsoft Azure (SQL DB, Data Factory), Git (GitHub, GitLab), Docker, CI/CD Pipelines \\
\textbf{Frontend \& UI/UX:} React, TypeScript, HTML5/CSS3, Responsive Design, User Experience Design \\
\textbf{Metodologías:} Scrum, Kanban, Documentación Técnica, Análisis Funcional, ETL Processes

\section{Formaci\'{o}n Acad\'{e}mica}
\textbf{Business Analytics Essentials, Machine Learning con Python, Data Science y Power BI}, Universidad de Chile – Escuela de Postgrado \hfill 2023 \\
\textbf{Optimización de Procesos y Administración de Bases de Datos SQL Server}, Universidad de Chile – Escuela de Postgrado \hfill 2023 \\
\textbf{Desarrollo de Apps Móviles}, Universidad Complutense de Madrid \hfill 2017 \\
\textbf{Analista Programador}, Universidad Tecnológica de Chile INACAP \hfill 2013

\end{document}